\nonstopmode
\documentclass[a4paper,11pt]{article}

\usepackage{fontspec}
\usepackage[a4paper, top=2.2cm, bottom=2.2cm, left=2.2cm, right=2.2cm]{geometry}
\usepackage{graphicx}
\usepackage{xcolor}
\usepackage{titlesec}
\usepackage{tabularx}
\usepackage{booktabs}
\usepackage{microtype}
\usepackage{hyperref}
\usepackage{enumitem}
\usepackage{array}
\usepackage{ragged2e}
\usepackage{parskip}
\usepackage{fancyhdr}
\usepackage{lastpage}
\usepackage{calc}
\usepackage{multirow}
\usepackage{mdframed}

\setmainfont[Ligatures=TeX]{Helvetica Neue}

\definecolor{raaRed}{HTML}{8B0000}
\definecolor{raaLight}{HTML}{f5f5dc}
\definecolor{raaMid}{HTML}{888888}
\definecolor{raaLine}{HTML}{cccccc}

\hypersetup{colorlinks=true, urlcolor=raaRed, linkcolor=raaRed}

\pagestyle{fancy}
\fancyhf{}
\renewcommand{\headrulewidth}{0pt}
\fancyfoot[L]{\footnotesize\color{raaMid} RAA.SPACE --- Matīsa iela 8, Rīga}
\fancyfoot[R]{\footnotesize\color{raaMid} \thepage\ / \pageref{LastPage}}

\titleformat{\section}{\large\bfseries\color{raaRed}}{}{0em}{}[\color{raaLine}\titlerule]
\titlespacing{\section}{0pt}{1.4em}{0.5em}

\titleformat{\subsection}{\normalsize\bfseries}{}{0em}{}
\titlespacing{\subsection}{0pt}{0.8em}{0.2em}

\newcommand{\evententry}[5]{%
  \noindent\begin{minipage}[t]{2.6cm}
    \vspace{0pt}
    \includegraphics[width=2.5cm,height=3cm,keepaspectratio]{#2}
  \end{minipage}%
  \hspace{0.3cm}%
  \begin{minipage}[t]{\linewidth - 2.9cm}
    \vspace{0pt}
    {\footnotesize\color{raaRed}\textbf{#1} \hfill \textit{\color{raaMid}#5}}\\[-0.1em]
    {\bfseries #3}\\[0.1em]
    {\small\RaggedRight #4}
  \end{minipage}%
  \par\vspace{0.4em}
  \noindent{\color{raaLine}\rule{\linewidth}{0.4pt}}
  \par\vspace{0.4em}
}

\newcommand{\milestone}[2]{%
  \noindent{\color{raaRed}\textbf{#1}} \quad #2\par\vspace{0.3em}
}

\newcommand{\teammember}[3]{%
  \begin{minipage}[t]{2.2cm}
    \vspace{0pt}
    \includegraphics[width=2cm,height=2.5cm,keepaspectratio]{#3}
  \end{minipage}%
  \hspace{0.3cm}%
  \begin{minipage}[t]{5.8cm}
    \vspace{0pt}
    {\bfseries #1}\\
    {\small\color{raaMid}#2}
  \end{minipage}%
}

\begin{document}

% ═════════════════════════════════════════════════════════
%  1. LAPPUSE — TITULLAPA + PAR RAA
% ═════════════════════════════════════════════════════════
\thispagestyle{empty}

\begin{center}
  \includegraphics[width=5cm,keepaspectratio]{/Users/leo/Projects/Raa/web/frontend/public/raa_logo_but_use_svg_if_possible.png}
  \par\vspace{0.6em}
  {\large\bfseries RAA.SPACE --- Kultūrtelpa}
  \par\vspace{0.2em}
  {\normalsize Radošā rūpnīca Veldze, Matīsa iela 8, Rīga LV-1001}
\end{center}

\vspace{0.4em}
\begin{center}
\includegraphics[width=\linewidth,height=8cm,keepaspectratio]{/Users/leo/Projects/Raa/web/frontend/public/assets/telpa/raa-galerija-1.jpg}
\end{center}
\vspace{0.4em}

\begin{center}
  \noindent{\color{raaLine}\rule{0.6\linewidth}{0.8pt}}
  \par\vspace{0.4em}
  \begin{tabular}{ccc}
    \textbf{E-pasts} & \textbf{Instagram} & \textbf{Adrese} \\[0.3em]
    \href{mailto:enter@raa.space}{enter@raa.space} &
    \href{https://instagram.com/raa_riga}{@raa\_riga} &
    Matīsa iela 8, Rīga LV-1001 \\
  \end{tabular}
  \par\vspace{0.4em}
  \noindent{\color{raaLine}\rule{0.6\linewidth}{0.8pt}}
  \par\vspace{0.3em}
  {\small\color{raaMid} Atsauces dokuments --- \today}
\end{center}

\vspace{0.5em}

% ─────────────────────────────────────────────────────────
\section{Par RAA}
% ─────────────────────────────────────────────────────────

RAA ir neatkarīga kultūrtelpa Radošajā rūpnīcā Veldze Rīgas centrā, Latvijā.
Tā atvērās \textbf{2025. gada 11. jūlijā}.

Telpa piedāvā daudzveidīgu programmu: laikmetīgās mākslas izstādes, teātra izrādes, koncertus,
dzejas lasījumus, darbnīcas, festivālus un kopienas pasākumus.

Telpa ietver izstāžu zāli (103\,m$^2$), ieejas vestibili (35\,m$^2$), teātra telpu
(13\,m$^2$) ar atsevišķu vestibili, aizkulišu koridoru un pieeju pagalmam, kas savieno
visu Veldzes kvartālu.

RAA (jeb \textit{reductio ad absurdum}) ved apgalvojumus līdz absurdam kopš 2025. gada.
RAA tev nav jābūt skaļam, lai tevi dzirdētu.

\vfill

\noindent{\color{raaLine}\rule{\linewidth}{0.4pt}}
{\footnotesize\color{raaMid}
  \textbf{RAA.SPACE} \enspace
  Biedrība \enspace
  Reģ.\ 40008349569 \enspace
  Matīsa iela 8, Rīga LV-1001
}

% ═════════════════════════════════════════════════════════
%  2. LAPPUSE — KOMANDA + MEŽA BĀRS + PLĀNOJUMS
% ═════════════════════════════════════════════════════════
\newpage

% ─────────────────────────────────────────────────────────
\section{Komanda}
% ─────────────────────────────────────────────────────────

\vspace{0.5em}
\noindent
\teammember{Leonhard Horstmeyer}{haosa meistars}{/Users/leo/Projects/Raa/web/frontend/public/assets/team/Leo.jpeg}%
\hspace{0.6cm}
\teammember{Jana Zariņa}{mākslinieciskā direktore}{/Users/leo/Projects/Raa/web/frontend/public/assets/team/JanaZarina.JPG}

\vspace{0.8em}
\noindent
\teammember{Helias Dzimoulakis}{teātra direktors}{/Users/leo/Projects/Raa/web/frontend/public/assets/team/HeliasD.jpg}

% ─────────────────────────────────────────────────────────
\section{MAATE NAATRE}
% ─────────────────────────────────────────────────────────

MAATE NAATRE ir meža bārs RAA telpas ieejas vestibilā, Radošajā rūpnīcā Veldze Rīgas centrā.
Atklāts 2026. gada 14. martā, tas ir neliela vieta, kur var aizkavēties, satikties ar citiem
un nogaršot dažādus dabas radītus dzērienus. Bārs atnes Latvijas mežu uz pilsētu ---
izmēģini dažādus kokteiļus no meža.

% ─────────────────────────────────────────────────────────
\section{Plānojums}
% ─────────────────────────────────────────────────────────

\vspace{0.5em}
\noindent
\begin{minipage}[t]{0.495\linewidth}
  \vspace{0pt}
  \includegraphics[width=\linewidth,height=7cm,keepaspectratio]{/Users/leo/Projects/Raa/web/frontend/public/assets/telpa/raa-galerija-0.jpg}
  \par\vspace{0.3em}
  {\small\color{raaMid}\textit{Galvenā zāle}}
\end{minipage}%
\hfill%
\begin{minipage}[t]{0.495\linewidth}
  \vspace{0pt}
  \includegraphics[width=\linewidth,height=7cm,keepaspectratio]{/Users/leo/Projects/Raa/web/frontend/public/assets/telpa/layout.jpeg}
  \par\vspace{0.3em}
  {\small\color{raaMid}\textit{RAA plānojums (bez sarkanās zonas)}}
\end{minipage}

% ═════════════════════════════════════════════════════════
%  3. LAPPUSE — HRONOLOĢIJA
% ═════════════════════════════════════════════════════════
\newpage

% ─────────────────────────────────────────────────────────
\section{Hronoloģija}
% ─────────────────────────────────────────────────────────

\milestone{Jūlijs 2025}{RAA atvērās 2025. gada 11. jūlijā. Atklāšanas grupu izstāde
  \textit{In Your Shoes}, kuratorēja Maija Kurševa, pulcēja septiņus māksliniekus no
  Latvijas Mākslas akadēmijas. Dzejas maratons \textit{Sinusoīda} ar Syntax Stupor.}

\milestone{Augusts 2025}{Katoptron Teātris, RAA rezidenta teātris, atvērās Helia Dzimoulaka
  vadībā. Arta Raituma vadīja divu dienu uguns darbnīcu, kam sekoja radīto darbu izstāde.
  Frīča Ripas soloizstāde \textit{Saules Gaismas Meklējumos}. Muzikāls vakars ar Elzu
  un Uģi Albiņiem. Ksenijas Kumačovas fotoizstāde.}

\milestone{Septembris 2025}{Homo Novus festivālam veltīta izstāde. Šujmašīnu tehno.
  Lielbritānijas dzejnieces Džoelas Teilores (T. S. Eliota balva) dzejas lasījums.
  Tvēruma \#11 numura prezentācija. Vakars ar Eriku Kolobanova un Konstantīnu Belihu.}

\milestone{Oktobris 2025}{Trīs dienu Starptautiskais smaržu mākslas festivāls (10.--12.\ okt.)
  ar māksliniekiem no dažādām valstīm. res Tacta performances un darbnīcu sērija.
  Brīvības Experience (16.--30.\ okt.) --- pilsētvides instalācija par Brīvības ielu,
  ko organizē Brunis Kibilds (Pilsēta cilvēkiem).}

\milestone{Novembris 2025}{Kefirs Studio grupu izstāde. Pirmais Rīgas Upcycling festivāls
  (15.\ nov.). Veronika Geidas fotoizstāde. \textit{Through the Looking Glass} ---
  Paulas Betijas Antenišķes fotogrāfija, dokumentējot Ivo Krieviņu un Rīgas ballroom scēnu.}

\milestone{Decembris 2025}{WunderKombināts publikācijas prezentācija. Klavierkoncerts ar
  Andreju Osokinsu. Tvēruma decembra numura prezentācija.}

\milestone{Janvāris 2026}{Ziemas pauze un renovācijas darbi RAA.}

\milestone{Februāris 2026}{30 dienu attīstoša instalācija \textit{AS DISTANT, SO NEAR} (29.\ janvāris -- 27.\ februāris) ar
  dažādiem māksliniekiem, kas pēta attāluma, tuvuma un klātbūtnes tēmas.}

\milestone{Marts 2026}{MAATE NAATRE meža bāra atklāšana (14.\ marts). Trīs pētniecības
  sanāksmju sesijas. Grāmatas prezentācija. Audēju grupu izstāde un Leonharda Horstmeyera
  performance \textit{Dialogo ad Infinitum}.}

\milestone{Aprīlis 2026}{\textit{ALTERNATĪVA: Pēdas no pazemes} --- 2ANNAS festivāla
  izstāde (10.\ apr. -- 4.\ maijs), ar četriem māksliniekiem un trim ekskursijām,
  ko atbalsta Rīgas dome un VKKF.}

\milestone{Maijs 2026}{LMA studentu grupu izstāde \textit{Atrodi mani, es esmu šeit}.
  Alīnas Burlakovas un Viktorijas Galkinas dueta izstāde \textit{TP054K-01xG}
  (15.\ maijs -- 8.\ jūnijs), tostarp Gallery Late vakars Riga Art Week ietvaros
  ar Elizabetes Balčus performanci.}

\milestone{Jūnijs 2026}{Rīgas Dizaina un mākslas vidusskolas ikgadējā produktu dizaina
  izstāde. Sākas \textit{Macchiato XVI --- Sauleszaķi}, starptautiskā kolektīva
  Macchiato Celebratio artis 16. festivāls.}

\milestone{Jūlijs 2026}{\textit{ITEM --- Shuffled Time} (5.--11.\ jūlijs) --- septiņu dienu
  performatīva nedēļa, kurā katrs mākslinieks saņem 24 stundu logu telpā.
  RAA pirmā gadadiena. Fotoklubs.lv fotoizstāde. \textit{RAA Šaha nedēļa}
  (20.--24.\ jūlijs) --- lekcijas, “Rokas un smadzenes” turnīrs, dzeja un teātris,
  kā arī šaha meistara F.-A.-D.\ Filidora baroka operas \textit{Blaise le Savetier}
  pirmatskaņojums Latvijā.}

\milestone{Augusts 2026}{\textit{Architecture of Feeling} --- lietuviešu ansambļa
  Sans Frayeur agrīnā baroka koncerts kopā ar viesmūziķiem no Itālijas un Latvijas.
  Eksperimentāla skaņas performance \textit{hg80}, Andra Alpa debijas dzejas krājuma
  \textit{vēl apmācies} atvēršana, Kristiāna Kļavas glezniecības soloizstāde
  \textit{Lidojums} un baroka mūzikas un leļļu teātra izrāde \textit{Dievs Pamet Johanu:
  Dzīve turpinās}.}

\milestone{Septembris 2026}{RAA starptautiskā jaunā teātra festivāla Homo Novus ietvaros
  (4.--12.\ septembris). Robert Rehnig skaņu instalāciju izstāde \textit{Collapsing Flat
  Waves} (24.\ septembris -- 14.\ oktobris).}

\milestone{Oktobris 2026}{Sākas \textit{RAA Ceturtdienas} --- iknedēļas ceturtdienu
  performanču sērija. Atkārtotas izrādes \textit{Dievs Pamet Johanu}. generation:Donut
  komunālā viesistaba Global Donut Days ietvaros. noscope.collective \textit{Kombināts 88}.
  Helmuta Bēķa romāna \textit{Cūcība Ķemeros} atklāšana ar grupu “Neīstā Atraitne”.}

\milestone{Novembris 2026}{Dāvja Vaivada stikla darbu izstāde \textit{Glass Boy}. Analogās
  fotogrāfijas izstāde \textit{Turās ainavā!} (16.\ novembris -- 8.\ decembris).
  \textit{Dievs Pamet Johanu} (27.--28.\ novembris).}

\milestone{Decembris 2026}{Džeza un stepa vakara džems (3.\ decembris). Koncerts
  \textit{Dowland: A Letter From Nuremberg} (4.\ decembris), atzīmējot 400 gadus kopš
  Džona Daulenda nāves.}

% ═════════════════════════════════════════════════════════
%  4. LAPPUSE+ — PASĀKUMI
% ═════════════════════════════════════════════════════════
\newpage

% ─────────────────────────────────────────────────────────
\section{Pasākumi}
% ─────────────────────────────────────────────────────────

{\small Visi RAA rīkotie pasākumi no atklāšanas 2025. gada 11. jūlijā hronoloģiskā secībā.
Pārējie 2026. gadam plānotie pasākumi tiks pievienoti, tiklīdz tie būs apstiprināti.}

\vspace{0.6em}

% ── JŪLIJS 2025 ──────────────────────────────────────────
\subsection*{\color{raaMid} Jūlijs 2025}

\evententry%
  {11. jūlijs 2025}%
  {/Users/leo/Projects/Raa/web/frontend/public/assets/raa-opening/banner-opening-thumb.jpeg}%
  {R A A Opening \#1}%
  {RAA kultūrtelpas atklāšana Matīsa ielā 8.}%
  {Atklāšana}

\evententry%
  {11.--18. jūlijs 2025}%
  {/Users/leo/Projects/Raa/web/frontend/public/assets/tavaas-kurpees/poster-thumb.jpg}%
  {In Your Shoes}%
  {Grupu izstāde, kuratorē Maija Kurševa (LMA KAS 3. kursa studenti).
  Katrs darbs atsaucas uz cita mākslinieka praksi. Septiņi mākslinieki.}%
  {Izstāde}

\evententry%
  {24.--27. jūlijs 2025}%
  {/Users/leo/Projects/Raa/web/frontend/public/assets/typewriter-museum/poster-thumb.jpg}%
  {Typewriter Museum Exhibition}%
  {20. gs. rakstāmmašīnu kolekcija no Vācijas, ASV, Zviedrijas un Itālijas.
  Apmeklētāji varēja rakstīt uz divām autentiskām mašīnām.}%
  {Izstāde}

\evententry%
  {27. jūlijs 2025}%
  {/Users/leo/Projects/Raa/web/frontend/public/assets/dzeja-sinusoiida/RAA_dzejasSinusidaPRINT-thumb.jpg}%
  {Dzeja: Sinusoīda}%
  {Dzejas maratons ar grupu Syntax Stupor, mūziku spēlē Dabsolis un Kaspars Sils.
  Atvērts mikrofons.}%
  {Dzeja / Mūzika}

% ── AUGUSTS 2025 ────────────────────────────────────────
\subsection*{\color{raaMid} Augusts 2025}

\evententry%
  {2.--3. augusts 2025}%
  {/Users/leo/Projects/Raa/web/frontend/public/assets/fire-workshop/uguns_poster_DRUKA_with_QR_code-thumb.png}%
  {Uguns darbnīca ar Artu Raitumu}%
  {Divu dienu eksperimentāla koka dedzināšanas darbnīca ar mākslinieci Artu Raitumu.}%
  {Darbnīca}

\evententry%
  {4.--6. augusts 2025}%
  {/Users/leo/Projects/Raa/web/frontend/public/assets/fire-workshop/uguns_poster_DRUKA_with_QR_code-thumb.png}%
  {Uguns darbnīcas izstāde}%
  {Darbnīcā radīto darbu izstāde.}%
  {Izstāde}

\evententry%
  {7.--10. augusts 2025}%
  {/Users/leo/Projects/Raa/web/frontend/public/assets/kas-aug-starp-mums/poster-thumb.png}%
  {Kas Aug Starp Mums}%
  {Ksenijas Kumačovas (dz. 1997) fotoizstāde.
  Pēta cilvēka un vietējās dabas mijiedarbību.}%
  {Izstāde}

\evententry%
  {9. augusts 2025}%
  {/Users/leo/Projects/Raa/web/frontend/public/assets/theatre-opening/poster-thumb.png}%
  {Katoptron Teātris atklāšana}%
  {RAA rezidenta teātra Katoptron Teātris atklāšana. Dibinājis un vada Helias Dzimoulakis.}%
  {Teātris}

\evententry%
  {14. augusts 2025}%
  {/Users/leo/Projects/Raa/web/frontend/public/assets/kabiirijas-salons/poster-thumb.jpg}%
  {Kabiirijas Salons \#1}%
  {Atvērta skatuve glezniecībai, mūzikai, performancei, dejai un dzejai.}%
  {Performance / Atvērtā skatuve}

\evententry%
  {14.--17. augusts 2025}%
  {/Users/leo/Projects/Raa/web/frontend/public/assets/shelf-stories/impression-thumb.jpeg}%
  {RAA Muzejs: Plauktu Stāsti}%
  {Muzejs ar priekšmetiem, ko kopienas locekļi ziedojuši ar saviem stāstiem.}%
  {Izstāde}

\evententry%
  {18.--26. augusts 2025}%
  {/Users/leo/Projects/Raa/web/frontend/public/assets/saules-gaismas-mekleejumos/poster1-thumb.jpeg}%
  {Saules Gaismas Meklējumos}%
  {Frīča Ripas soloizstāde. Atklāšanā DJ seti.}%
  {Izstāde}

\evententry%
  {28. augusts 2025}%
  {/Users/leo/Projects/Raa/web/frontend/public/assets/laime-pliist-uz-traukiem/poster-thumb.jpg}%
  {Laime Plīst uz Traukiem}%
  {Vakara uzveduma pasākums.}%
  {Performance}

% ── SEPTEMBRIS 2025 ─────────────────────────────────────
\subsection*{\color{raaMid} Septembris 2025}

\evententry%
  {1. septembris 2025}%
  {/Users/leo/Projects/Raa/web/frontend/public/assets/second-variety-show/poster-thumb.jpg}%
  {Second Variety Show}%
  {Dažādu priekšnesumu vakars.}%
  {Performance}

\evententry%
  {3. septembris 2025}%
  {/Users/leo/Projects/Raa/web/frontend/public/assets/pilseetas-naktii/poster-thumb.jpg}%
  {Pilsētas Naktī}%
  {Vakara pasākums par pilsētvides tematiku.}%
  {Pasākums}

\evententry%
  {8. septembris 2025}%
  {/Users/leo/Projects/Raa/web/frontend/public/assets/dzoela-teilore/Dzoela_Teilore_Raa_A3-thumb.jpg}%
  {Džoela Teilore --- dzejas lasījums}%
  {Dzejas lasījums ar Lielbritānijas dzejnieci Džoelu Teilori, T. S. Eliota balvas laureāti.}%
  {Literatūra / Dzeja}

\evententry%
  {14. septembris 2025}%
  {/Users/leo/Projects/Raa/web/frontend/public/assets/sujmasinu-tehno/poster-thumb.png}%
  {Šujmašīnu Tehno}%
  {Vakars, kas apvieno amatniecību, kopienu un ritmu.}%
  {Darbnīca / Pasākums}

\evententry%
  {17. septembris 2025}%
  {/Users/leo/Projects/Raa/web/frontend/public/assets/tautas-dziesmas/tautas-poster-thumb.png}%
  {Tautas dziesmas}%
  {Radoša pēcpusdiena un vakars ar dziedāšanu, rakstīšanu, zīmēšanu un dalīšanos ---
  atvērta telpa, kurā iepazīt tautas dziesmas caur mūziku un radošu izpausmi.}%
  {Mūzika / Darbnīca}

\evententry%
  {18. septembris 2025}%
  {/Users/leo/Projects/Raa/web/frontend/public/assets/ja-nu-saule-rit-neuzaust/poster-thumb.jpeg}%
  {Tvērums \#11 ``Austrumi'': Ko darīt, ja saule rīt neuzaust}%
  {Literārā žurnāla Tvērums 11. numura prezentācija.}%
  {Literatūra / Prezentācija}

\evententry%
  {27. septembris 2025}%
  {/Users/leo/Projects/Raa/web/frontend/public/assets/exo-foto-izstade/poster-thumb.jpeg}%
  {Exo --- fotoizstāde}%
  {Fotoizstāde.}%
  {Izstāde}

% ── OKTOBRIS 2025 ───────────────────────────────────────
\subsection*{\color{raaMid} Oktobris 2025}

\evententry%
  {4. oktobris 2025}%
  {/Users/leo/Projects/Raa/web/frontend/public/assets/kizomba/poster-thumb.png}%
  {Kizomba / UrbanKiz ballīte}%
  {Sociālās dejas ballīte ar Kizomba un UrbanKiz mūziku.}%
  {Deja / Ballīte}

\evententry%
  {5. oktobris 2025}%
  {/Users/leo/Projects/Raa/web/frontend/public/assets/riga-sock-exchange/poster-thumb.png}%
  {Rīgas zeķu maiņa}%
  {Kopienas apģērbu maiņas pasākums.}%
  {Kopiena}

\evententry%
  {8. oktobris 2025}%
  {/Users/leo/Projects/Raa/web/frontend/public/assets/res-tacta/poster-thumb.png}%
  {res tacta --- performance / darbnīca}%
  {Performance un darbnīca, kas pēta taustes pieredzi.}%
  {Performance / Darbnīca}

\evententry%
  {10.--12. oktobris 2025}%
  {/Users/leo/Projects/Raa/web/frontend/public/assets/scent-art-festival/poster-thumb.jpg}%
  {Starptautiskais smaržu mākslas festivāls}%
  {Trīs dienu starptautisks festivāls, kas veltīts ožas mākslinieciskajai izpausmei ---
  izstādes, performances, darbnīcas un smaržu ekskursija pa Rīgu.}%
  {Festivāls}

\evententry%
  {16.--30. oktobris 2025}%
  {/Users/leo/Projects/Raa/web/frontend/public/assets/brivibas-experience/poster-thumb.png}%
  {Brīvības Experience --- pilsētvides instalācija}%
  {Izstāde un pilsētvides instalācija, kas veltīta Rīgas galvenajai ielai, Brīvības ielai.
  Izmantojot vizualizācijas, 3D modeļus un citus medijus, projekts pēta ielas vēsturi,
  pašreizējo stāvokli un nākotnes scenārijus. Organizators: Brunis Kibilds
  (Pilsēta cilvēkiem). Apmeklētāji tika aicināti piedalīties diskusijā par
  cilvēkdraudzīgāku pilsētvidi.}%
  {Pasākums}

% ── NOVEMBRIS 2025 ──────────────────────────────────────
\subsection*{\color{raaMid} Novembris 2025}

\evententry%
  {6.--9. novembris 2025}%
  {/Users/leo/Projects/Raa/web/frontend/public/assets/kefirs-studio-exhibition/poster-thumb.jpg}%
  {Kefirs Studio izstāde}%
  {Kefirs Studio grupu izstāde.}%
  {Izstāde}

\evententry%
  {15. novembris 2025}%
  {/Users/leo/Projects/Raa/web/frontend/public/assets/pirmais-rigas-upcycling-festivals/poster-thumb.jpg}%
  {Pirmais Rīgas Upcycling festivāls}%
  {Publisks pasākums, kas veicina ilgtspējīgu dizainu un otrreizēju pārstrādi.
  Darbnīcas, tirgus un kopienas aktivitātes.}%
  {Festivāls}

\evententry%
  {17. novembris 2025}%
  {/Users/leo/Projects/Raa/web/frontend/public/assets/pa-liekas-izstade/poster-thumb.jpeg}%
  {(pa)liekas --- fotoizstāde}%
  {Veronika Geidas soloizstāde.}%
  {Izstāde}

\evententry%
  {24.--30. novembris 2025}%
  {/Users/leo/Projects/Raa/web/frontend/public/assets/through-the-looking-glass/poster-thumb.jpg}%
  {Through the Looking Glass}%
  {Fotoizstāde, kas pēta kvīru identitāti un ballroom kultūru Latvijā.
  Fotogrāfija --- Paula Betija Antenišķe, dokumentējot Ivo Krieviņu un Rīgas ballroom scēnu.
  Atbalstīts ar Eiropas Savienības līdzfinansējumu.}%
  {Izstāde}

% ── DECEMBRIS 2025 ──────────────────────────────────────
\subsection*{\color{raaMid} Decembris 2025}

\evententry%
  {9. decembris 2025}%
  {/Users/leo/Projects/Raa/web/frontend/public/assets/wunderkombinats/poster-thumb.jpg}%
  {WunderKombināts --- prezentācija}%
  {WunderKombināts publikācijas prezentācija.}%
  {Literatūra / Prezentācija}

\evententry%
  {12. decembris 2025}%
  {/Users/leo/Projects/Raa/web/frontend/public/assets/andrejs-osokins-koncerts/poster-thumb.jpg}%
  {Andrejs Osokins --- klavierkoncerts}%
  {Klavierkoncerts ar Andreju Osokinsu.}%
  {Koncerts}

\evententry%
  {19. decembris 2025}%
  {/Users/leo/Projects/Raa/web/frontend/public/assets/tverums-decembris/poster-thumb.jpg}%
  {Tvērums --- decembra numura prezentācija}%
  {Tvēruma decembra numura prezentācija.}%
  {Literatūra / Prezentācija}

\evententry%
  {28. decembris 2025}%
  {/Users/leo/Projects/Raa/web/frontend/public/assets/marks-ozolins-koncerts/poster-thumb.jpg}%
  {Mārks Ozoliņš --- koncetvakars}%
  {Koncetvakars ar Mārku Ozoliņu.}%
  {Koncerts}

% ── JANVĀRIS 2026 ───────────────────────────────────────
\subsection*{\color{raaMid} Janvāris 2026}

\evententry%
  {29. janvāris -- 27. februāris 2026}%
  {/Users/leo/Projects/Raa/web/frontend/public/assets/as-far-so-near/poster-thumb.jpeg}%
  {AS DISTANT, SO NEAR --- instalācija}%
  {30 dienu attīstoša instalācija ar dažādiem atklātajā konkursā atlasītiem māksliniekiem,
  kas pēta attāluma, tuvuma un klātbūtnes tēmas.}%
  {Instalācija / Atklātais konkurss}

% ── MARTS 2026 ─────────────────────────────────────────
\subsection*{\color{raaMid} Marts 2026}

\evententry%
  {14. marts 2026}%
  {/Users/leo/Projects/Raa/web/frontend/public/assets/upcycling-tirgus-march/poster-thumb.jpg}%
  {2. Upcycling festivāls ``Atjauno veco!''}%
  {Otrais Rīgas Upcycling festivāls. Tirgus, darbnīcas un kopienas aktivitātes.}%
  {Festivāls}

\evententry%
  {14. marts 2026}%
  {/Users/leo/Projects/Raa/web/frontend/public/assets/maate-naatre-atklaashana/poster-thumb.jpg}%
  {MAATE NAATRE --- atklāšana}%
  {MAATE NAATRE meža bāra atklāšana RAA telpā.}%
  {Atklāšana}

\evententry%
  {17., 24.\ un 27. marts 2026}%
  {/Users/leo/Projects/Raa/web/frontend/public/assets/citizen-research-meetup-first-meeting/poster-thumb.png}%
  {Pētniecības sanāksmju sērija}%
  {Trīs atvērtas sanāksmes RAA zināšanu apmaiņai un kolektīvai izpētei.
  \textbf{Pilsoņu pētniecības sanāksme} (17.\ marts) --- atvērtu sesiju sākums datu
  jautājumu apspriešanai un jaunu prasmju apguvei. Organizē jopalesi.science.
  \textbf{Rekreatīvās pētniecības sanāksme} (24.\ marts) --- neformāla ideju
  un zinātkāru jautājumu apmaiņa.
  \textbf{Ārkārtas sanāksme: Rīmana dzeta} (27.\ marts) --- īpaša sesija par
  Rīmana dzeta funkciju un matemātisko skaistumu.}%
  {Sanāksme / Pētniecība}

\evententry%
  {26. marts -- 6. aprīlis 2026}%
  {/Users/leo/Projects/Raa/web/frontend/public/assets/weavers/poster-thumb.png}%
  {Audēji --- izstāde}%
  {Tekstila un starpnozaru izstāde.}%
  {Izstāde}

\evententry%
  {26. marts 2026}%
  {/Users/leo/Projects/Raa/web/frontend/public/assets/weavers/poster-thumb.png}%
  {Dialogo ad Infinitum --- performance}%
  {Leonharda Horstmeyera performance. Daļa no Audēju izstādes.}%
  {Performance}

\evententry%
  {28. marts 2026}%
  {/Users/leo/Projects/Raa/web/frontend/public/assets/laimes-govs/poster-thumb.jpg}%
  {Laimes Govs}%
  {Kopienas rīta pasākums.}%
  {Pasākums}

% ── APRĪLIS -- MAIJS 2026 ──────────────────────────────
\subsection*{\color{raaMid} Aprīlis -- Maijs 2026}

\evententry%
  {10. aprīlis -- 4. maijs 2026}%
  {/Users/leo/Projects/Raa/web/frontend/public/assets/alternativa-2-annas/poster-thumb.jpeg}%
  {ALTERNATĪVA: Pēdas no pazemes}%
  {2ANNAS festivāla izstāde. Mākslinieki: Laimdota Malle (LV), alliance Kurss (LV),
  Rebecca Korn (UK), Anastasija Šnepe-Šnepe (LV). Kurators: Kristiāna Kārkliņa.
  Trīs ekskursijas. Atbalsta Rīgas dome un VKKF.}%
  {Izstāde}

\evententry%
  {8.--10. maijs 2026}%
  {/Users/leo/Projects/Raa/web/frontend/public/assets/es-esmu-seit/poster-thumb.jpg}%
  {Atrodi mani, es esmu šeit}%
  {LMA studentu grupu izstāde kursa ``Izstāžu veidošanas prakse māksliniekiem'' ietvaros.
  Kuratorē Ieva Stalšene.}%
  {Izstāde}

\evententry%
  {15. maijs -- 8. jūnijs 2026}%
  {/Users/leo/Projects/Raa/web/frontend/public/assets/TP054K-01xG/poster-thumb.jpeg}%
  {TP054K-01xG}%
  {Dueta izstāde --- Alīna Burlakova (dz. 1992) un Viktorija Galkina (dz. 1988).
  Atklāšana 15. maijā ar Eleonoras Kampes vokālo performance.}%
  {Izstāde}

\evententry%
  {28. maijs 2026}%
  {/Users/leo/Projects/Raa/web/frontend/public/assets/TP054K-01xG/gallery-late-poster-thumb.jpeg}%
  {Gallery Late (RAW): TP054K-01xG}%
  {Gallery Late vakars Riga Art Week ietvaros izstādē TP054K-01xG
  ar Elizabetes Balčus performanci.}%
  {Performance / Gallery Late}

% ── JŪNIJS -- AUGUSTS 2026 ─────────────────────────────
\subsection*{\color{raaMid} Jūnijs -- Augusts 2026}

\evententry%
  {12.--22. jūnijs 2026}%
  {/Users/leo/Projects/Raa/web/frontend/public/assets/produktu-dizaina-izstaade/poster-thumb.jpeg}%
  {Produktu dizaina izstāde 2026}%
  {Rīgas Dizaina un mākslas vidusskolas (RDMV) ikgadējā produktu dizaina izstāde,
  kurā skatāms audzēkņu gada radošais darbs ar tēmu \textit{Sezons no Slēgums}.
  Atklāšana 11. jūnijā.}%
  {Izstāde}

\evententry%
  {26. jūnijs -- 3. jūlijs 2026}%
  {/Users/leo/Projects/Raa/web/frontend/public/assets/Macchiato/poster-macchiato-thumb.jpeg}%
  {Macchiato XVI --- Sauleszaķi}%
  {Starptautiskā kolektīva Macchiato Celebratio artis 16. festivāls. Nedēļu ilga
  multimediāla izstāde ar septiņiem māksliniekiem, kā arī koncerti, literatūras vakars,
  laikmetīgā deja, performances, darbnīcas, algorave un DJ setu programma.}%
  {Festivāls}

\evententry%
  {5.--11. jūlijs 2026}%
  {/Users/leo/Projects/Raa/web/frontend/public/assets/item/public-poster-thumb.png}%
  {ITEM --- Shuffled Time}%
  {Septiņu dienu performatīva nedēļa. Katru dienu viens mākslinieks vai mākslinieku grupa
  saņem 24 stundu logu, kurā Invoke, Temper, Elate, Mint telpu; katrā dienā notiek arī
  iepriekš plānota intervence.}%
  {Festivāls / Performance}

\evententry%
  {13.--18. jūlijs 2026}%
  {/Users/leo/Projects/Raa/web/frontend/public/assets/fotoklubs-photo-exhibition/poster-thumb.png}%
  {Fotoklubs.lv fotoizstāde}%
  {Fotoklubs.lv fotoizstāde, atvērta katru dienu no plkst. 15.00 līdz 20.00.}%
  {Izstāde}

\evententry%
  {20.--24. jūlijs 2026}%
  {/Users/leo/Projects/Raa/web/frontend/public/assets/raa-chess-week/poster-thumb.png}%
  {RAA Šaha nedēļa}%
  {Nedēļa, kas veltīta šaham un tā kultūras nospiedumam; sadarbībā ar Ex Libris katru dienu
  neformālā gaisotnē tika spēlēts šahs. Atklāta 20. jūlijā --- Starptautiskajā šaha dienā,
  tās 60. gadā --- ar Bruno Ranča un Signes Kacēnas lekcijām, un noslēgta 24. jūlijā
  ar Santas Didžus un Leonharda Horstmeijera teatrālām šaha ainām.}%
  {Festivāls}

\evententry%
  {21. jūlijs 2026}%
  {/Users/leo/Projects/Raa/web/frontend/public/assets/hands-and-brain-tournament/poster-thumb.png}%
  {Rokas un smadzenes turnīrs}%
  {Šaha turnīrs pāros --- viens spēlētājs pārvieto figūras, otrs nosauc tikai to, kura
  figūra jāpārvieto. Organizē Signe Kacēna un Visvaldis Gercāns kopā ar RAA un Ex Libris.
  RAA Šaha nedēļas ietvaros.}%
  {Turnīrs}

\evententry%
  {22. jūlijs 2026}%
  {/Users/leo/Projects/Raa/web/frontend/public/assets/blaise-le-savetier/poster-thumb.png}%
  {Blaise le Savetier}%
  {Komiskās viencēliena baroka operas pirmatskaņojums Rīgā. Autors --- komponists un
  leģendārais šahists Fransuā Andrē Danikans Filidors. Koncertizrādi sniedz Lietuvas
  baroka ansamblis \textit{Lieto Fine} kopā ar latviešu instrumentālistiem;
  režisore Evelina Brėdikytė. RAA Šaha nedēļas ietvaros.}%
  {Opera}

\evententry%
  {23. jūlijs 2026}%
  {/Users/leo/Projects/Raa/web/frontend/public/assets/vienkarsi-sarezgiti/poster-thumb.png}%
  {Vienkārši sarežģīti}%
  {Kolektīva “Sintakses stupors” deviņu dzejnieku lasījumi un Zoya Zoyroko dziesmas.
  RAA Šaha nedēļas ietvaros.}%
  {Dzeja / Mūzika}

\evententry%
  {5. augusts 2026}%
  {/Users/leo/Projects/Raa/web/frontend/public/assets/architecture-of-feeling/poster-thumb.jpeg}%
  {Architecture of Feeling}%
  {Agrīnā baroka koncerts --- ansamblis Sans Frayeur (LT): Justina Kaminskaitė
  (mecosoprāns) un Skirmantė Kirkliauskaitė (klavesīns), kopā ar viesiem
  Marco Baronchelli (lauta, IT), Marianna Puriņa (baroka vijole) un
  Elfrīda Švalbe (baroka čells).}%
  {Koncerts}

\evententry%
  {6. augusts 2026}%
  {/Users/leo/Projects/Raa/web/frontend/public/assets/hg80/poster-thumb.jpeg}%
  {hg80 --- LAIKLEJUPSLĪDE}%
  {Eksperimentāla skaņas performance par dzīvsudraba tēmu --- diviem čelliem, balsij un basklarnetei.}%
  {Performance / Mūzika}

\evententry%
  {7. augusts 2026}%
  {/Users/leo/Projects/Raa/web/frontend/public/assets/vel-apmacies/poster-thumb.jpeg}%
  {vēl apmācies}%
  {Andra Alpa debijas dzejas krājuma \textit{vēl apmācies} (Orbīta) atvēršana.}%
  {Grāmatas atvēršana / Dzeja}

\evententry%
  {14.--22. augusts 2026}%
  {/Users/leo/Projects/Raa/web/frontend/public/assets/lidojums/poster-thumb.jpeg}%
  {Lidojums --- Kristiāns Kļava}%
  {Kristiāna Kļavas glezniecības soloizstāde par trauksmi, karu un pārmaiņām.}%
  {Izstāde}

\evententry%
  {29. aug., 16. un 18. okt., 27. un 28. nov. 2026}%
  {/Users/leo/Projects/Raa/web/frontend/public/assets/dievs-pamet-johanu/poster-thumb.jpeg}%
  {Dievs Pamet Johanu: Dzīve turpinās}%
  {Baroka mūzikas vizuālā teātra izrāde --- baroka mūzikas un leļļu teātra saspēle par 17.\ gadsimta vijolnieka Johana Fišera ceļu cauri Jelgavai un Rīgai. Lomās Santa Didžus un Guna Zariņa, senās mūzikas apvienība Concerto Riga; muzikālā vadītāja Laura Šarova, režija Helias Dzimoulakis. Atbalsta VKKF.}%
  {Teātris / Mūzika}

% ── SEPTEMBRIS -- OKTOBRIS 2026 ────────────────────────
\subsection*{\color{raaMid} Septembris -- Oktobris 2026}

\evententry%
  {4.--12. septembris 2026}%
  {/Users/leo/Projects/Raa/web/frontend/public/assets/homo-novus-2026/poster-thumb.jpeg}%
  {Homo Novus 2026}%
  {Starptautiskais jaunā teātra festivāls Homo Novus ar notikumiem Daugavas krastos, Rīgā un bijušās Veldzes fabrikas radošajā kvartālā, kur mājo RAA.}%
  {Festivāls}

\evententry%
  {24. septembris -- 14. oktobris 2026}%
  {/Users/leo/Projects/Raa/web/frontend/public/assets/collapsing-flat-waves/poster-2-thumb.jpeg}%
  {Collapsing Flat Waves --- Robert Rehnig}%
  {Leipcigas skaņu mākslinieka Robert Rehnig skaņu instalāciju izstāde par telpu, klusumu un pāreju no dzirdēšanas uz klausīšanos.}%
  {Izstāde / Skaņa}

\evententry%
  {1. oktobris -- 3. decembris 2026}%
  {/Users/leo/Projects/Raa/web/frontend/public/assets/raa-ceturtdienas-okt-01/main-thumb.jpg}%
  {RAA Ceturtdienas}%
  {Iknedēļas ceturtdienu sērija, kurā RAA kļūst par skatuvi īsām izrādēm, skatuves eksperimentiem, lasījumiem un miniatūrām. \textbf{1.\ okt.} --- Teātra studija “Mēs” un Annija Sprīvule. \textbf{8.\ okt.} --- Katerinas “cardigan performance” un Elīnas Vendijas Rībenas “Bēdiņš un Sauliņš”. \textbf{22.\ okt.} --- Anonīma “Pick Up The Pieces” un Ernesta Ā “Nr.\ 5”. \textbf{5.\ nov.} --- “Vēstule draudzenei”, Patrīcijas M.\ Keišas un Marijas Luīzes Meļķes lasījumi. \textbf{12.\ nov.} --- “[mõ:tus] Skaņa kustībā” un “THE NIGHTMARE HAS OPENING HOURS”, pēc tam dzīvā mūzika. \textbf{3.\ dec.} --- džeza un stepa vakara džems ar Žaneti Muižnieci, Anniju Sustrupi un Eduardu Hermani.}%
  {Performanču sērija}

\evententry%
  {23. oktobris 2026}%
  {/Users/leo/Projects/Raa/web/frontend/public/assets/gen-donut-viesistaba/poster-thumb.jpg}%
  {gen:Donut komunālā viesistaba \#1}%
  {Global Donut Days ietvaros --- pirmā generation:Donut tīklojuma kopradītā komunālā viesistaba: ceļojums uz 2036, Rīgas ilgtspējas (virtuļa) portrets, mājas DJ sets un kopīga nākamo sanākšanu plānošana. Līdzfinansē Eiropas Savienība (GearUP!, ES Klimata pakts).}%
  {Kopiena}

\evententry%
  {24. oktobris 2026}%
  {/Users/leo/Projects/Raa/web/frontend/public/assets/kombinats-88/poster-thumb.jpg}%
  {Kombināts 88 --- noscope.collective}%
  {Uz vienu dienu noscope.collective pārņēma Veldzi --- tetovēšana, pīrsings, māksla, darbnīcas un mūzika no 13:00 līdz pusnaktij.}%
  {Festivāls / Māksla}

\evententry%
  {29. oktobris 2026}%
  {/Users/leo/Projects/Raa/web/frontend/public/assets/cuciba-kemeros/poster-thumb.jpeg}%
  {Cūcība Ķemeros --- grāmatas atklāšana}%
  {Helmuta Bēķa kriminālromāna \textit{Cūcība Ķemeros} atklāšanas svētki, kuros grupa “Neīstā Atraitne” pirmo reizi dzīvajā izpildīja trīs romānam radītus skaņdarbus, tostarp debijas singlu “Melu taure”.}%
  {Grāmatas atklāšana / Koncerts}

% ── NOVEMBRIS -- DECEMBRIS 2026 ────────────────────────
\subsection*{\color{raaMid} Novembris -- Decembris 2026}

\evententry%
  {7.--14. novembris 2026}%
  {/Users/leo/Projects/Raa/web/frontend/public/assets/glassboy/glassboy-poster-thumb.jpeg}%
  {Glass Boy --- Dāvis Vaivads}%
  {Dāvja Vaivada stikla darbu izstāde: darbi no 2021.\ gada līdz mūsdienām un jauni darbi. Atklāšanā performance un DJ Kibernācija vinila sets; muzikāls vakars 12.\ novembrī ar MEETRA, Jātnieks bez rumpja un DAFELZ.}%
  {Izstāde}

\evententry%
  {16. novembris -- 8. decembris 2026}%
  {/Users/leo/Projects/Raa/web/frontend/public/assets/turas-ainava/main-poster-thumb.jpg}%
  {Turās ainavā!}%
  {Mārica Kraukles, Katrīnas Lilientāles un Emīla Stuķēna-Viļumsona analogās fotogrāfijas izstāde --- Andreja Granta foto meistardarbnīca RTJN “Annas 2”.}%
  {Izstāde}

\evententry%
  {4. decembris 2026}%
  {/Users/leo/Projects/Raa/web/frontend/public/assets/dowland-a-letter-from-nuremberg/poster-thumb.jpeg}%
  {Dowland: A Letter From Nuremberg}%
  {Koncerts, atzīmējot 400 gadus kopš Džona Daulenda nāves: renesanses dziesmas un instrumentālā mūzika ar lasījumiem Early Modern English valodā no viņa 1595.\ gada vēstules seram Robertam Sesilam. Karolina Parmas (soprāns, LT), Laura Šarova (vijole, LV), Jana Zariņa (klavesīns, LV), Leonhard Horstmeyer (aktieris, DE).}%
  {Koncerts}

\end{document}
